\chapter{Bayes 解析の基礎}\label{ch:17}

\begin{goalbox}
\begin{itemize}
\item 事前分布・尤度・事後分布・事前予測分布（周辺分布）・事後予測分布を区別して書ける．
\item Beta--Binomial の共役事前分布による事後分布・事後平均・事後モードを導出できる．
\item 事後平均が「MLE と事前平均の重み付き平均」（縮小）になることを説明できる．
\item 一様事前分布の下で，$Y$ の周辺分布が離散一様分布になることを示し，事後確率 $\Prob(\pi\ge\pi_0\mid y)$ による判定ができる．
\item 人時間を考慮したポアソンモデルで，率の最尤推定量 $\hat\lambda=Y/T$ を求められる．
\end{itemize}
\end{goalbox}

\begin{exambox}
\begin{itemize}
\item \kakomon{2022}{4}：有効例数 $Y\sim\Bin(n,\pi)$，一様事前分布．尤度と MLE，事後分布 $\Be(y+1,n-y+1)$，事後モード $=y/n$，
      事後平均 $(y+1)/(n+2)$ と MLE・事前平均の関係，$Y$ の周辺分布（離散一様），$\Prob(\pi\ge1/2\mid y)\ge0.8$ となる最小の $y$（$n=5$ で $y=4$）．
\item \kakomon{2025}{2}〔1〕：人口の異なる2市の死亡数（ポアソン）から共通の死亡率の MLE（1.8）．
      同問〔2〕〜〔5〕（率の差の検定，ガンマ事前分布と事後予測分布）は，分散の推定量や事前密度が問題文で与えられる誘導付きの設問である（\cref{app:F}）．
\item 統計数理では\kakomonSM{2021}{3}（ポアソン平均の十分性・MLE・スコア型区間），\kakomonSM{2024}{3}（Beta 事前分布の事後平均と同じ形の推定量）．
\end{itemize}
\end{exambox}

\flowline{二項データ\\件数＋人時間}{比率・率}{共役 Bayes\\ポアソン MLE}{事前分布\\二項・ポアソン性}{事後 $=$ 事前 $\times$ 尤度}{縮小と\\事前の影響}

\section{ポアソン率の最尤推定}\label{sec:17-poisson}
観察人時間（または人口）$T$ の集団で起きたイベント数を $Y$ とし，率 $\lambda$（単位人時間あたり）で
\[
Y\sim\Po(\lambda T),\qquad \Prob(Y=y)=\frac{e^{-\lambda T}(\lambda T)^y}{y!}
\]
とする．\kakomon{2025}{2} のように「1万人あたり $\Po(\theta)$」が独立なら，ポアソン分布の再生性により人口 $T$ 万人の死亡数は $\Po(T\theta)$ である．
\begin{meidai}\label[meidai]{prop:17-mle}
$\hat\lambda=Y/T$，$\V[\hat\lambda]=\lambda/T$（推定値 $Y/T^2$），デルタ法により $\V[\log\hat\lambda]\approx1/(\lambda T)$（推定値 $1/Y$）．
複数の集団で率が共通なら $\hat\lambda=\sum_g Y_g/\sum_g T_g$．
\end{meidai}
\begin{proof}
$\log L=y\log\lambda-\lambda T+\text{const}$，$\partial\log L/\partial\lambda=y/\lambda-T=0$．共通の率では $\log L=(\sum_g y_g)\log\lambda-\lambda\sum_g T_g$ で同様．
$\V[Y]=\lambda T$ より $\V[\hat\lambda]=\lambda/T$．$\log$ の微分 $1/\lambda$ を掛けて $\V[\log\hat\lambda]\approx(1/\lambda)^2\lambda/T$．
\end{proof}
\kakomon{2025}{2}〔1〕：$\hat\theta=(4+14)/(2+8)=1.8$（1万人あたり）．
率の近似95\%信頼区間は $\exp\bigl(\log\hat\lambda\pm1.96/\sqrt Y\bigr)$（$z_{0.025}=1.96$）である．
この尤度 $\lambda^ye^{-\lambda T}$ は，打ち切りのある指数モデルの尤度 $\lambda^de^{-\lambda W}$（\cref{prop:10-expmle}）と同じ形である．

\section{Bayes 推測の基本}\label{sec:17-bayes}
\Hissu
パラメータ $\theta$ に事前分布 $p(\theta)$ を置き，データ $y$ を観測した後の事後分布を
\[
p(\theta\mid y)=\frac{p(y\mid\theta)p(\theta)}{p(y)}\propto p(y\mid\theta)\,p(\theta),\qquad p(y)=\int p(y\mid\theta)p(\theta)\,\dd\theta
\]
とする．$p(y)$ は\emphx{事前予測分布}（周辺分布）であり，将来のデータ $\tilde y$ の\emphx{事後予測分布}は
\[
p(\tilde y\mid y)=\int p(\tilde y\mid\theta)\,p(\theta\mid y)\,\dd\theta
\]
（$\tilde y$ と $y$ が $\theta$ を与えて独立な場合）．
点推定には事後平均（2乗損失の Bayes 推定量），事後モード（事後分布の最頻値），事後中央値を，区間推定には事後分布の $2.5\%$ 点と $97.5\%$ 点による信用区間を用いる．
事後分布が事前分布と同じ分布族になる事前分布を\emphx{共役事前分布}という．

\section{Beta--Binomial モデル}\label{sec:17-betabin}
\Hissu
\begin{teiri}[Beta--Binomial]\label[teiri]{thm:17-betabin}
$Y\mid\pi\sim\Bin(n,\pi)$，$\pi\sim\Be(a,b)$ なら
\begin{equation}
\pi\mid y\sim\Be(a+y,\ b+n-y),\qquad \E[\pi\mid y]=\frac{a+y}{a+b+n}=\frac{n}{a+b+n}\cdot\frac yn+\frac{a+b}{a+b+n}\cdot\frac{a}{a+b}.\label{eq:17-betabin}
\end{equation}
事後モードは $(a+y-1)/(a+b+n-2)$（$a+y\ge1$，$b+n-y\ge1$，$a+b+n>2$ のとき）．$Y$ の周辺分布は
$\Prob(Y=y)=\binom ny\dfrac{B(a+y,b+n-y)}{B(a,b)}$（ベータ二項分布）．
\end{teiri}
\begin{proof}
$p(\pi\mid y)\propto\pi^y(1-\pi)^{n-y}\cdot\pi^{a-1}(1-\pi)^{b-1}=\pi^{a+y-1}(1-\pi)^{b+n-y-1}$ はベータ分布の核．
平均は $\Be(\alpha,\beta)$ の平均 $\alpha/(\alpha+\beta)$，モードは密度の対数を微分して0とおいた $(\alpha-1)/(\alpha+\beta-2)$．
周辺分布は $\int_0^1\binom ny\pi^y(1-\pi)^{n-y}\frac{\pi^{a-1}(1-\pi)^{b-1}}{B(a,b)}\dd\pi$ をベータ関数で表したもの．
\end{proof}
事後平均は MLE $y/n$ と事前平均 $a/(a+b)$ の重み付き平均で，重みは「データの大きさ $n$」と「事前分布の大きさ $a+b$（事前の仮想的な標本サイズ）」の比である．
$n\to\infty$ で事前の影響は消える．

\paragraph{一様事前分布（\kakomon{2022}{4}）}
$a=b=1$ なら事後分布 $\Be(y+1,n-y+1)$，事後モード $y/n$（MLE と一致），事後平均 $\frac{y+1}{n+2}=\frac{n}{n+2}\cdot\frac yn+\frac{2}{n+2}\cdot\frac12$．
$Y$ の周辺分布は
\[
\Prob(Y=y)=\binom ny B(y+1,n-y+1)=\frac{n!}{y!(n-y)!}\cdot\frac{y!(n-y)!}{(n+1)!}=\frac{1}{n+1}\quad(y=0,\dots,n),
\]
すなわち離散一様分布である．

\paragraph{事後確率による判定}
「$\Prob(\pi\ge\pi_0\mid y)\ge\gamma$ なら有効」とする判定では，事後確率は
$\Prob(\pi\ge\pi_0\mid y)=\dfrac{(n+1)!}{y!(n-y)!}\displaystyle\int_{\pi_0}^1\pi^y(1-\pi)^{n-y}\dd\pi$（一様事前分布）で，$y$ について増加する．
積分の代わりに次の関係を使うこともできる．
\begin{meidai}\label[meidai]{prop:17-betatail}
$\pi\mid y\sim\Be(y+1,n-y+1)$ のとき $\Prob(\pi>p\mid y)=\Prob(W\le y)$，$W\sim\Bin(n+1,p)$．
\end{meidai}
\begin{proof}
$m$ 個の独立な $\Unif(0,1)$ の $k$ 番目に小さい値 $U_{(k)}$ は $\Be(k,m-k+1)$ に従い，
「$U_{(k)}\le p$」は「$m$ 個のうち $k$ 個以上が $p$ 以下」と同値なので $\Prob(\Be(k,m-k+1)\le p)=\Prob(\Bin(m,p)\ge k)$．
これを $k=y+1$，$m=n+1$ で用い，余事象をとる．
\end{proof}
\kakomon{2022}{4} では $n=5$，$p=1/2$ で $\Prob(\pi\ge\frac12\mid y=3)=\Prob(\Bin(6,\tfrac12)\le3)=42/64=0.656$，
$y=4$ で $57/64=0.891$．よって $0.8$ 以上となる最小の $y$ は4．
（問題文で与えられる積分値を使っても，$y=3$ で $60\times\frac{7}{640}=0.656$，$y=4$ で $30\times\frac{19}{640}=0.891$ と同じ値になる．）

\paragraph{事後予測}
次の $\tilde n$ 人での有効数 $\tilde Y$ の事後予測分布は，$\Be(a+y,b+n-y)$ を事前分布とみなしたベータ二項分布で，
次の1人が有効である確率は事後平均 $\E[\pi\mid y]$ に等しい．

\section{正規--正規モデル}\label{sec:17-normal}
$\bar X\mid\mu\sim\Normal(\mu,\sigma^2/n)$（$\sigma^2$ 既知），$\mu\sim\Normal(\xi,\tau^2)$ なら
\[
\mu\mid\bar x\sim\Normal\Bigl(\frac{(n/\sigma^2)\bar x+(1/\tau^2)\xi}{n/\sigma^2+1/\tau^2},\ \frac{1}{n/\sigma^2+1/\tau^2}\Bigr).
\]
事後平均は精度（分散の逆数）で重み付けた $\bar x$ と事前平均 $\xi$ の平均であり，Beta--Binomial と同じく縮小の形をしている．

\section{頻度論との違い}\label{sec:17-vs}
\begin{itemize}
\item 頻度論ではパラメータは定数で，確率は手順の繰り返しに関するもの．Bayes ではパラメータの不確かさを確率分布で表し，「$\Prob(\pi\ge0.5\mid\text{データ})$」を直接述べられる．
\item 信用区間は「データを与えたとき $\theta$ が区間に入る確率が95\%」と解釈できる．信頼区間はそうは解釈できない．
\item 事前分布が結論に影響する．小標本では影響が大きく，事前分布の選択の根拠と感度分析が必要である．
\item Bayes の判定規則（事後確率の閾値）も，頻度論的な性質（第1種の過誤確率・検出力）を評価できる（\cref{reidai:17-2}(2)）．
\end{itemize}

\section{仮定}\label{sec:17-assump}
\begin{itemize}
\item 二項モデル：各人の結果が独立で，有効確率 $\pi$ が共通．
\item ポアソンモデル：イベントが独立で，率が観察期間・集団内で一定．
\item Bayes：事前分布の妥当性．共役事前分布は計算の便宜であり，それ自体が正しい保証はない．
\end{itemize}

\section{結果の解釈}\label{sec:17-interp}
\begin{itemize}
\item 「事後平均 0.30，95\%信用区間 0.15--0.47」：データと事前情報を合わせると，有効率は95\%の確率でこの範囲にある．
\item 事後平均が MLE と違うのは，事前平均への縮小による．データが多いほど両者は近づく．
\item 一様事前分布では事後モードは MLE と一致するが，事後平均は $1/2$ の側に寄る．
\end{itemize}

\section{手法選択}\label{sec:17-choice}
\begin{itemize}
\item 小標本の2値アウトカムで既存情報（または一様事前）を使いたい → Beta--Binomial の事後分布と事後確率による判定．
\item 将来の患者の有効数の予測 → 事後予測分布（ベータ二項）．
\item 観察期間・人口が異なる件数データの率 → ポアソンモデルの MLE $\hat\lambda=Y/T$．
\end{itemize}

\section{よくある誤り}\label{sec:17-err}
\begin{warnbox}
\begin{itemize}
\item 事後モードと事後平均を混同する（一様事前で MLE と一致するのはモード）．
\item ベータ分布のパラメータを $\Be(y,n-y)$ とする（一様事前 $\Be(1,1)$ の1を足し忘れる）．
\item 事前予測分布（周辺分布）と事後予測分布を混同する．
\item 事後予測を，事後平均を代入した二項分布 $\Bin(\tilde n,\E[\pi\mid y])$ とする（$\pi$ の不確かさを無視）．
\item 人時間を無視して件数どうしを比べる．
\end{itemize}
\end{warnbox}

\section{数理統計との接続}\label{sec:17-link}
\begin{linkbox}
\begin{itemize}
\item ベータ--二項（『現代数理統計学の基礎』例4.14，例6.12）．
\item 二項の累積確率とベータ分布の関係（同書第3章演習問19）が\cref{prop:17-betatail}である．
\item \kakomonSM{2024}{3} の MSE 一定の推定量 $(S_n+\sqrt n/2)/(n+\sqrt n)$ は $\Be(\sqrt n/2,\sqrt n/2)$ 事前分布の事後平均と同じ形である．
\item \kakomonSM{2021}{3} のポアソン平均の区間 $\frac1n(t+2\mp2\sqrt{t+1})$ は，正規近似 $(t-n\lambda)^2\le4n\lambda$ を解いたスコア型区間である．
\end{itemize}
\end{linkbox}

\section{例題}\label{sec:17-ex}
\begin{reidai}[Beta--Binomial]\label{reidai:17-1}
先行研究から有効率の事前分布を $\Be(2,8)$ とし，新しい試験で20人中7人が有効であった．
\begin{enumerate}[label=(\arabic*)]
\item 事前平均，事後分布，事後平均，事後モードを求めよ．
\item 事後平均を MLE と事前平均の重み付き平均として表せ．
\item （任意）次の2人の有効数 $\tilde Y$ の事後予測分布を求めよ（2人の結果は $\pi$ を与えて互いに独立で，既存データとも独立とする）．
\end{enumerate}
\end{reidai}

\begin{reidai}[事後確率による判定]\label{reidai:17-2}
一様事前分布の下，$n=9$ 人の試験で「$\Prob(\pi>0.5\mid y)\ge0.9$ なら有望」と判定する．
$\Bin(10,0.5)$ の累積確率 $\Prob(W\le y)$ は $y=5,6,7,8$ でそれぞれ $0.623,0.828,0.945,0.989$ である．
\begin{enumerate}[label=(\arabic*)]
\item 有望と判定される最小の有効数を求めよ．
\item 真の有効率が0.5のとき，有望と判定される確率を求めよ（$\Bin(9,0.5)$ で $\Prob(Y\ge7)=0.0898$）．
\end{enumerate}
\end{reidai}

\section{解答}\label{sec:17-sol}
\begin{kaitou}{reidai:17-1}
(1) 事前平均 $2/10=0.2$．事後 $\Be(2+7,\,8+13)=\Be(9,21)$，平均 $9/30=0.30$，モード $8/28=0.286$．\\
(2) $0.30=\frac{20}{30}\times0.35+\frac{10}{30}\times0.2$（事前の仮想標本10人，データ20人，MLE $7/20=0.35$）．\\
(3)（任意）$\Prob(\tilde Y=0\mid y)=\E[(1-\pi)^2\mid y]=\frac{21}{30}\cdot\frac{22}{31}=\frac{462}{930}=0.497$，
$\Prob(\tilde Y=2\mid y)=\E[\pi^2\mid y]=\frac{9}{30}\cdot\frac{10}{31}=\frac{90}{930}=0.097$，
$\Prob(\tilde Y=1\mid y)=2\E[\pi(1-\pi)\mid y]=2\cdot\frac{9\cdot21}{30\cdot31}=\frac{378}{930}=0.406$．
（$\Be(\alpha,\beta)$ で $\E[\pi^2]=\frac{\alpha(\alpha+1)}{(\alpha+\beta)(\alpha+\beta+1)}$ など．）
\end{kaitou}

\begin{kaitou}{reidai:17-2}
(1) \cref{prop:17-betatail}より $\Prob(\pi>0.5\mid y)=\Prob(W\le y)$，$W\sim\Bin(10,0.5)$．$\ge0.9$ となる最小の $y$ は7（$0.945$）．\\
(2) $\Prob(Y\ge7\mid\pi=0.5)=0.0898$．この Bayes 判定の「第1種の過誤確率」は約9\%である．
\end{kaitou}

\begin{summarybox}
\begin{itemize}
\item 事後 $\propto$ 尤度 $\times$ 事前．事前予測分布 $p(y)=\int p(y\mid\theta)p(\theta)\dd\theta$，事後予測分布 $p(\tilde y\mid y)=\int p(\tilde y\mid\theta)p(\theta\mid y)\dd\theta$．
\item Beta--Binomial：$\Be(a+y,b+n-y)$，事後平均は MLE と事前平均の重み付き平均（重み $n$ と $a+b$）．
\item 一様事前：事後 $\Be(y+1,n-y+1)$，モード $y/n$，平均 $(y+1)/(n+2)$，$Y$ の周辺分布は離散一様 $1/(n+1)$．
\item 事後確率による判定：$\Prob(\pi>p\mid y)=\Prob(\Bin(n+1,p)\le y)$．
\item 正規--正規：事後平均は精度で重み付けた平均．
\item ポアソン率：$Y\sim\Po(\lambda T)$ で $\hat\lambda=Y/T$，$\V[\log\hat\lambda]\approx1/Y$．
\end{itemize}
\end{summarybox}
